\documentclass[12pt]{article}
\usepackage[utf8]{inputenc}
\usepackage[T1]{fontenc}
\usepackage[margin=1in]{geometry}
\usepackage{setspace}
\usepackage{xcolor}
\usepackage[colorlinks=true, urlcolor=blue!60!black, linkcolor=blue!60!black]{hyperref}
\setlength{\parindent}{0em}
\setlength{\parskip}{0.6em}
\renewcommand{\baselinestretch}{1.0}
\begin{document}
\begin{center}
{\LARGE \textit{Yick Wo v. Hopkins}, 118 U.S. 356 (1886)\par}
\vspace{0.5em}
{\large Case Brief --- Two-Page Edition\par}
\end{center}

\section{Memory Jogger}
A unanimous Court held that San Francisco's facially neutral laundry ordinance, administered to deny permits to every Chinese applicant, violated equal protection, and that the Fourteenth Amendment protects noncitizens as ``persons'' (\href{https://supreme.justia.com/cases/federal/us/118/356/}{118 U.S. 356}).

\section{Detailed Case Facts}
San Francisco required the Board of Supervisors' consent to operate a laundry in a wooden building, with no criteria for consent. Yick Wo, a Chinese subject who had run a safe, inspected laundry for twenty-two years, was refused without reason. All roughly 200 Chinese applicants were denied, while all but one of about eighty non-Chinese applicants were approved.

\section{Procedural History}
Yick Wo was convicted, fined \$10, and jailed. The California Supreme Court denied habeas; in the companion \textit{Wo Lee} case, a federal circuit judge found the ordinance aimed at Chinese laundries but deferred. The Supreme Court heard both cases.

\section{Judicial Votes}
Reversed, 9--0. Justice Matthews wrote for a unanimous Court.

\section{Holding}
Equal protection extends to all persons within a State's jurisdiction, regardless of nationality. A neutral law administered to discriminate between similarly situated persons by race or nationality violates the Fourteenth Amendment, as does a grant of naked, arbitrary power over a lawful occupation.

\section{Analysis of Opinions}
Matthews read the ordinances as conferring ``a naked and arbitrary power'' over persons as well as places, unlike the uniform curfew upheld in \textit{Barbier}. The Fourteenth Amendment's shift from ``citizens'' to ``any person'' protected Chinese residents. Holding one's livelihood at another's will is ``the essence of slavery itself.'' Given the admitted disparity, no explanation existed but hostility to the petitioners' race and nationality. \textit{Washington v. Davis} later recast \textit{Yick Wo} as a rare case of disparity so stark that purpose was inescapable.

\section{Examples: Future Applications}
\textbf{Same-Side:} a standardless permit board denying 142 of 145 immigrant applicants with clean inspections while approving others. \textbf{Opposite-Side:} a uniform, mechanically enforced fire code with disparate impact but no purpose is valid under \textit{Davis}. \textbf{Fence-Sitter:} a 45\% versus 20\% denial gap with mixed explanations and one hostile remark.

\section{Critique}
The opinion blends an objection to arbitrary power with an anti-discrimination principle, and later cases kept only the second. Its stipulated record makes it an evidentiary outlier, so treating it as the benchmark makes subtler discrimination hard to prove.

\section{Key Quotations}
\begin{quote}
``Though the law itself be fair on its face and impartial in appearance, yet, if it is applied and administered by public authority with an evil eye and an unequal hand, so as practically to make unjust and illegal discriminations between persons in similar circumstances, material to their rights, the denial of equal justice is still within the prohibition of the Constitution.''
\end{quote}
\noindent --- Justice Matthews, Opinion of the Court, 118 U.S. at 373--374.

\end{document}
